\par\Needspace{6\baselineskip}\medskip
\hypertarget{native-112-article083}{}\pdfbookmark[2]{83.}{article-083}
\begin{center}83.\end{center}
If we prefer to begin solving the same problem by calculating the angle $\Pi$, we use the following method. Of equations I we subtract the third from the second, the third from the first, and the second from the first. This gives the following three new equations.
\[
\begin{aligned}
&\frac{\dfrac1{r'}-\dfrac1{r''}}{2\sin\tfrac12(N''-N')}=\frac{e}{p}\sin(\tfrac12N'+\mathpunct{\cdot}\tfrac12N''-\Pi)\\[6pt]
\mathrm{II.}\qquad&\frac{\dfrac1r-\dfrac1{r''}}{2\sin\tfrac12(N''-N)}=\frac{e}{p}\sin(\tfrac12N+\tfrac12N''-\Pi)\\[6pt]
&\frac{\dfrac1r-\dfrac1{r'}}{2\sin\tfrac12(N'-N)}=\frac{e}{p}\sin(\tfrac12N+\tfrac12N'-\Pi).
\end{aligned}
\]
Any two of these equations give $\Pi$ and $\dfrac{e}{p}$ by Lemma II of Article 78. From these values we also obtain $e$ and $p$ using any one of equations I. If we adopt the third solution given in Article 78, II, combining the first and third of equations II gives the following algorithm. Determine the auxiliary angle $\zeta$ from the equation
\[
\operatorname{tang}\zeta=\frac{\dfrac{r'}r-1}{1-\dfrac{r'}{r''}}\cdot\frac{\sin\tfrac12(N''-N')}{\sin\tfrac12(N'-N)}\mathpunct{\raisebox{.45ex}{,}}
\]
\SourceLeaf{113}and then we have
\[
\operatorname{tang}(\tfrac14N+\tfrac12N'+\tfrac14N''-\Pi)=\operatorname{tang}(45^0+\zeta)\operatorname{tang}\tfrac14(N''-N).
\]
Interchanging the second position with either the first or the third gives two other entirely analogous solutions. Since this method gives less convenient formulas for $\dfrac{e}{p}$, it is preferable to obtain $e$ and $p$ from two of equations I by the method of Article 80. The ambiguity in determining $\Pi$ from the tangent of the angle $\tfrac14N+\tfrac12N'+\tfrac14N''-\Pi$ must be resolved so that $e$ is positive. The values obtained for $e$ are clearly opposite when the values chosen for $\Pi$ differ by $180^0$. The sign of $p$, however, is independent of this ambiguity. The value of $p$ can be negative only if the three given points lie on the part of the hyperbola turned away from the Sun. This case is contrary to the laws of nature and lies outside our discussion here.

What would result from applying the first method of Article 78, II, after more laborious substitutions, can in the present case be obtained more conveniently as follows. Multiply the first of equations II by $\cos\tfrac12(N''-N')$ and the third by $\cos\tfrac12(N'-N)$, then subtract the latter product from the former. Applying Lemma I of Article 78 appropriately\footnote{Specifically, in the second formula set $A=\tfrac12(N''-N')$, $B=\tfrac12N+\tfrac12N''-\Pi$, $C=\tfrac12(N-N'')$.} then yields the equation
\[
\begin{aligned}
&\tfrac12\left(\dfrac1{r'}-\dfrac1{r''}\right)\operatorname{cotang}\tfrac12(N''-N')-\tfrac12\left(\dfrac1r-\dfrac1{r'}\right)\operatorname{cotang}\tfrac12(N'-N)\\[4pt]
&\qquad=\frac{e}{p}\sin\tfrac12(N''-N)\cos(\tfrac12N+\tfrac12N''-\Pi).
\end{aligned}
\]
Combining this equation with the second of equations II gives $\Pi$ and $\dfrac{e}{p}$, with $\Pi$ given by the formula
\[
\begin{aligned}
&\operatorname{tang}(\tfrac12N+\tfrac12N''-\Pi)\\[4pt]
&\qquad=\frac{\dfrac{r'}r-\dfrac{r'}{r''}}{\left(1-\dfrac{r'}{r''}\right)\operatorname{cotang}\tfrac12(N''-N')-\left(\dfrac{r'}r-1\right)\operatorname{cotang}\tfrac12(N'-N)}\mathpunct{\cdot}
\end{aligned}
\]
Two other entirely analogous formulas can also be derived from this one by interchanging the second position with either the first or the third.
